%============================================================
\chapter{Ergebnisse und Evaluation}
\label{chap:ergebnisse}
%============================================================

\section{Testumgebung}
\label{sec:test_env}

Alle Versuchsläufe wurden auf einem Indoor-Parcours mit den
folgenden Eigenschaften durchgeführt:

\begin{itemize}
  \item Abmessungen: ca.\ $6\,\text{m} \times 4\,\text{m}$
  \item Bodenbelag: Linoleumboden (reflektierend)
  \item Hindernisse: 4--6 statische Pylone (Höhe ca.\ $50\,\text{cm}$)
  \item Beleuchtung: Deckenleuchten (kein Tages\-lichteintrag)
\end{itemize}

Die Tests wurden in zwei Phasen unterteilt:
\begin{enumerate}
  \item \textbf{Simulator} (F1TENTH Gym): Erste Funktionsverifikation.
  \item \textbf{Reales Fahrzeug}: Übertragung auf die Hardware-Plattform.
\end{enumerate}

%------------------------------------------------------------
\section{Simulationsergebnisse}
\label{sec:sim_results}
%------------------------------------------------------------

Im Simulator konnte der Algorithmus bei Verwendung der in
\cref{tab:parameter} angegebenen Standardparameter in allen
Testläufen mindestens \textbf{80\,\%} der Hindernisse erfolgreich
umfahren.
Bei geringer Hindernisdichte (2--3 Pylone auf dem Parcours)
wurde ein Durchfahren ohne Kollision in \(8\) von \(10\) Läufen
beobachtet.

\begin{table}[H]
  \centering
  \caption{Simulationsergebnisse (F1TENTH Gym)}
  \label{tab:sim_results}
  \begin{tabular}{@{}lcccc@{}}
    \toprule
    \textbf{Szenario}  & \textbf{Läufe}
      & \textbf{Kollisionsfrei} & \textbf{Kollisionen}
      & \textbf{Erfolgsrate} \\
    \midrule
    Wenige Hindernisse (2--3)  & 10 & 8 & 2 & 80\,\% \\
    Mittlere Dichte (4--5)     & 10 & 5 & 5 & 50\,\% \\
    Hohe Dichte (6+)           & 10 & 2 & 8 & 20\,\% \\
    \bottomrule
  \end{tabular}
\end{table}

%------------------------------------------------------------
\section{Ergebnisse auf dem realen Fahrzeug}
\label{sec:real_results}
%------------------------------------------------------------

Bei der Übertragung auf das reale Fahrzeug traten erhebliche
Leistungseinbrüche auf.
In den meisten Testläufen kollidierte das Fahrzeug oder verlor
die Kontrolle innerhalb der ersten $5\,\text{s}$ nach dem Start.
Die beobachteten Hauptprobleme werden im folgenden Abschnitt
dokumentiert.

\begin{table}[H]
  \centering
  \caption{Ergebnisse auf dem realen Fahrzeug}
  \label{tab:real_results}
  \begin{tabular}{@{}lcccc@{}}
    \toprule
    \textbf{Szenario}  & \textbf{Läufe}
      & \textbf{Kollisionsfrei} & \textbf{Kollisionen}
      & \textbf{Erfolgsrate} \\
    \midrule
    Leerer Parcours               & 5 & 4 & 1 & 80\,\% \\
    Wenige Hindernisse (2--3)     & 5 & 1 & 4 & 20\,\% \\
    Mittlere Dichte (4--5)        & 5 & 0 & 5 &  0\,\% \\
    \bottomrule
  \end{tabular}
\end{table}

%------------------------------------------------------------
\section{Identifizierte Probleme}
\label{sec:problems}
%------------------------------------------------------------

\subsection{LiDAR-Reflexionen am Boden}
\label{subsec:reflections}

Der reflektierende Linoleumboden erzeugte Spiegelungspunkte,
die vom LiDAR als nahe Hindernisse interpretiert wurden.
Dies führte dazu, dass der Algorithmus fiktive Sicherheitsblasen
erzeugte und freie Lücken fälschlicherweise versperrte.

\textbf{Mögliche Abhilfemaßnahme:} Filterpipeline zur Entfernung
von Boden-Reflexionen (z.\,B.\ Höhen\-filter via IMU-Fusion).

\subsection{Sim-to-Real-Gap}
\label{subsec:sim2real}

Die im Simulator erzielten Parameter ließen sich auf das reale
Fahrzeug nicht unmittelbar übertragen.
Die Motorcharakteristik, das Reifenverhalten sowie Latenzzeiten
zwischen Sensor und Aktor wichen von der Simulation ab.

\textbf{Mögliche Abhilfemaßnahme:} Systematisches Param-Tuning
direkt auf der Hardware oder Identifikation eines genaueren
Fahrzeugmodells für die Simulation.

\subsection{Zu hohe Fahrgeschwindigkeit}
\label{subsec:speed}

Bei $v_{\max} = 1{,}5\,\text{m/s}$ reagierte das Fahrzeug
bei unvorhergesehenen Hindernissen zu träge.
Eine Reduktion auf $v_{\max} = 0{,}8\,\text{m/s}$ verbesserte
das Ausweichverhalten, erhöhte aber die Gesamtfahrtzeit deutlich.

\subsection{Unzureichendes Blasen-Tuning}
\label{subsec:bubble_tuning}

Der Blasenradius $r_{\text{bubble}} = 0{,}35\,\text{m}$ war
für die realen Hindernisabmessungen zu groß gewählt und verdeckte
bei engen Durchfahrten die einzige freie Lücke vollständig.
